\documentclass{article}
\usepackage[hidelinks]{hyperref}
\title{A durable effect does not release a remote hold until its result arrives}
\begin{document}
\maketitle

\paragraph{Integration claim.}
The \href{https://github.com/eternal-november/eos-research-programmes/pull/631}{guarded pledge}
requires qualified terminal evidence. The
\href{https://github.com/eternal-november/eos-research-programmes/pull/633}{durable warehouse adapter}
provides an authoritative local result after a worker crash. The
combined simulation joins them through actual sends and callbacks of
the pinned SubT broker. The protected source and the consuming station
are two distinct native warehouse instances. The source stays off its
named bay through the consuming station's one required native step
and should return by the final available confirmation cut if it can.
An independent source carrier can still move and send at ticks 2 and 6.
The cross-plant meaning of the duty, capability, clock ordering and
route authorization are defined by the wrapper, not by the native
programs. No equipment or independent mission is used.

\paragraph{Evidence and obligations.}
The station worker exits silently in one of four phases around
publication and is actually relaunched in the positive arm of every
joint case. The operation either committed at cut 5 or still needs its
one application at cut 6. A paid receipt originates only from the
durable result for the exact application key, whose native first-stage
result is checked. A receipt payload reaching the broker callback still
does not constitute evidence for another key. The fixture transmits
an actual three-byte value, with a provisioned namespace and trusted
sender paths; it does not implement authentication, replay-resistant
identity allocation or an external durability protocol.

Denote by $E$ the authoritative station effect, by $R$ the first
qualifying callback, and by $B$ source reentry. For a positive receipt
these cuts satisfy $E\le R\le B$, with the callback ordered first if
reentry shares an outer tick. The local return guard checks the actual
proposed native step before executing it. Lost evidence may therefore
prevent return even when a durable effect exists. Cancellation is
outside this fixture: absence of the paid receipt is not terminal
unpaid proof. No cross-world snapshot or shared simulator transaction
is used to release; only the delivered named result is consumed.

\paragraph{Radio boundary and replay.}
The first confirmation opportunity at tick 6 uses P and spends three
of its five bytes even if RF loses the packet. Its transmit debit is
still live at tick 7; an erased attempt does not spend the receiver's
capacity. A preauthorized Q can therefore carry confirmation at tick 7
if that link works. If neither real confirmation arrives, the guard
remains. A delayed old but valid receipt for a different key can be
delivered at tick 1; accepting it is wrong even though its packet and
terminal status are genuine for its own operation.

Across 64 configurations per arm, the qualified adapter returns in 40
and preserves source clearance at the station effect in all. A direct
radio check using the same information and the same guarded release
ties, so this is not a composition speedup. A primary-only route returns
in 32 safely; it cannot use the added preauthorized backup, so it is a
route-capability ablation rather than a matched state-of-the-art stack.
Trusting a planned commit breaks clearance in 48 configurations;
accepting the unrelated terminal result for this key breaks it in 32.
All independent movements and the actual independent callbacks run.
The aggressive comparator arms use the same verified station cut trace
generated in the positive execution, not a modified station effect
authority. Counts describe the finite design sweep, not failure rates.

\paragraph{Remaining boundary.}
This joins effect-local simulated recovery with remotely acquired,
resource-priced evidence. It does not make a powered external effect
rollback-capable, recover a source guard across a whole-host crash,
or promise liveness under permanent receipt loss. Stronger disruption
must be matched by the actual capabilities and budget of the competing
system. Independent transfer, whole-stack inventory dispositions and
full I1 remain open.
\end{document}
